\documentclass[11pt]{article}

\usepackage[utf8]{inputenc}
\usepackage[T1]{fontenc}
\usepackage{amsmath}
\usepackage{amssymb}
\usepackage{amsthm}
\usepackage{mathtools}
\usepackage[hidelinks]{hyperref}
\usepackage{doi}
\hypersetup{
  unicode,
  pdftitle={Chiral Orientation of the Directed Heisenberg Frontier},
  pdfauthor={Jérôme Beau}
}

\theoremstyle{plain}
\newtheorem{theorem}{Theorem}
\newtheorem{proposition}{Proposition}
\newtheorem{lemma}{Lemma}
\theoremstyle{remark}
\newtheorem{remark}{Remark}
\newtheorem*{status}{Status}

\newcommand{\Zq}{\mathbb{Z}/q\mathbb{Z}}
\newcommand{\zq}{\zeta_q}
\newcommand{\dac}{\Delta A_c}
\newcommand{\rhochi}{\rho_\chi}
\newcommand{\thmax}{\theta_{\max}}
\newcommand{\JPi}{J_\Pi}
\newcommand{\Rmix}{R_{\mathrm{mix}}}
\newcommand{\NA}{N_A}
\newcommand{\EPi}{E_\Pi}
\newcommand{\PiS}{\Pi_S}
\newcommand{\sgn}{\operatorname{sign}}
\newcommand{\two}{\mathbf{2}}
\newcommand{\twobar}{\bar{\mathbf{2}}}
\newcommand{\Horient}{[\mathrm{H\text{-}orient}]}

\title{Chiral Orientation of the Directed Heisenberg Frontier}
\author{J\'er\^ome Beau\\
Independent Researcher\\
\texttt{jerome.beau@cosmochrony.org}}
\date{2026}

\begin{document}
  \maketitle

  \begin{abstract}
    The directed Heisenberg frontier carries a chirally dressed cocycle signal
    $\Theta_\chi = \langle \sigma_L\,(2\pi/q)\,\dac\rangle_{\partial^+}$, a discrete counterpart of the orientation
    datum (in the sense of~\cite{Beau2026q11of}) of the angular generation split of the projected Dirac residue
    $\EPi$; it is not the $J_3$ coefficient of the metaplectic step, which in the $\mathrm{SL}(2)$ model of Q14 is
    produced by the $\mathfrak{sl}_2$ commutator $[E,F]=H$ (Q14, Proposition~6.5).
    The raw cocycle average and the maximal-locking bound $\thmax$ are computable on the Heisenberg frontier of~\cite{Beau2026q11of}, but the
    signal itself depends on the chiral weight $\sigma_L$, whose origin lies in the Lorentzian spin-lift and not in the
    arithmetic frontier.
    We prove, under the hypothesis [H-WS] of~\cite{Beau2026q11of} (the residual reflection $R_b$ acts on the chiral
    carrier as a scalar, so $[\gamma_5,R_b]=0$, which is all that this note uses; the value $-1$ of the scalar is a second part,
used only in~\cite{Beau2026aog}; supplied by no source), that $\sigma_L$ cannot be any function of the real cascade generator $T(e)$: edge inversion and the residual
    reflection $\varphi$ both act as $T\mapsto -T$, so every $f(T)$ has identical arrow-parity and $\varphi$-parity,
    whereas $\sigma_L$ must be arrow-even and $\varphi$-odd.
    The resolution, under [H-WS], is that $\sigma_L$ is the Weyl representation type, $\two$ versus $\twobar$: group
    inversion preserves the representation (arrow-even) while complex conjugation exchanges it ($\varphi$-odd).
    In the inherited Heisenberg central-phase lift the chiral phase is the linear cyclotomic character
    $\omega_\chi(e)=\zq^{\dac(e)}$, with no Gauss, Maslov, or $\sqrt{5}$ component; on the nonzero-cocycle support the
    natural branch satisfies $\sigma_L(e)=\sgn \Im\,\zq^{\dac(e)} = \sgn\dac(e)$, giving, by this definition, $\rhochi = 1$ and
    $\NA = \thmax \neq 0$.
    This establishes $\Horient$ in the inherited Heisenberg central-phase lift and under [H-WS], not unconditionally:
    the only remaining obstruction besides [H-WS] is an additional spin-Galois phase outside the central character, in particular a $\sqrt{5}$ or
    ADE-spin factor orthogonal to $\zq$.
    The result is a conditional derivation of the frontier amplitude datum; it does not derive $u\neq 0$, which would
    further require a transport of the oriented signal to the $\JPi$-odd generation sector; the only transport
    stated in~\cite{Beau2026prs} is the hypothesis [H-Tr] for the $\mathfrak{sl}_2$ opening, supplied by no source,
    and the identification of the frontier signal with that opening is a further input.
  \end{abstract}

  \noindent\textbf{Keywords:} directed Heisenberg frontier; chiral orientation; Weil representation type; metaplectic
  central character; generation split; Schur complement; arithmetic orthogonality; epistemic stratification.

  \newpage
  \tableofcontents

  \newpage
  \section{Introduction}

    The fermionic sub-programme locates the inter-generation mass split in the $J_3$-odd part of $\EPi^2$ on the
    gauge-singlet generation factor $\mathbb{C}^3_{\mathrm{gen}} \cong \mathrm{Sym}^2(V)$, read through the reading
    map $\Phi$ of the hypothesis [H-Gen] of~\cite{Beau2026prs}, which is supplied there as a hypothesis.
    That note shows the Schur form $-N^\dagger N\preceq0$ of the compression remainder of the projected Dirac square,
    which is the projective endomorphism of Q14 only under the identity $\PiS D^2\PiS^*=\mathrm{Lich}$, and derives no
    mechanism for a non-zero split.
    The oriented-frontier diagnostic of the recursive $Q11$ model defines a single observable on the directed Cayley
    graph of $\mathrm{Heis}_3(\Zq)$~\cite{Beau2026q11of}, a discrete counterpart of the orientation datum (in the sense of~\cite{Beau2026q11of}) of this
    split; no map from it to the $J_3$-odd part of $\EPi^2$ is supplied.
    That diagnostic returns two computable quantities and one missing datum.
    The raw directed cocycle average vanishes,
    \begin{equation}
      \langle \dac \rangle_{\partial^+} = 0,
    \end{equation}
    because the residual reflection $\varphi$ pairs the outgoing frontier with itself and sends $\dac\mapsto -\dac$.
    The maximal-locking bound,
    \begin{equation}
      \thmax = \Big\langle \tfrac{2\pi}{q}\,|\dac| \Big\rangle_{\partial^+},
    \end{equation}
    is also computable, and $q$-stable (see below).
    The signal that would turn the bound into an amplitude,
    \begin{equation}
      \Theta_\chi = \Big\langle \sigma_L\,\tfrac{2\pi}{q}\,\dac \Big\rangle_{\partial^+},
      \qquad |\Theta_\chi| \le \thmax,
    \end{equation}
    requires the chiral weight $\sigma_L(e)\in\{\pm 1\}$.
    The frontier computation does not supply $\sigma_L$: it is the surviving stub of the Lorentzian chiral lift, and the
    question $\Horient$ is whether $\sigma_L$ is non-orthogonal to $\sgn\dac$ on the outgoing frontier,
    \begin{equation}
      \rhochi = \frac{\langle \sigma_L\,\dac\rangle_{\partial^+}}{\langle |\dac|\rangle_{\partial^+}}\neq 0.
    \end{equation}

    Both quantities are exact: with the generators $\{\pm X,\pm Y\}$ of~\cite{Beau2026a16}, the raw average and the
    symmetrised anti-bias control vanish by the $\varphi$-pairing above, and on each shell $\thmax\cdot q$ is $2\pi$
    times a rational number independent of $q$ for $q\ge 61$ ($\tfrac13,\tfrac59,\tfrac23,\ldots$ for
    $m=1,2,3,\ldots$; the script of~\cite{Beau2026q11of}, which weights shells by the capacity increment, gives
    $\thmax\cdot q=2.114,\ 2.122,\ 2.118$ at $n_3^{\mathrm{obs}}=2$ for $q=61,101,151$, stable to $0.4\,\%$, not
    invariant), as recomputed by \texttt{code/frontier\_exact.py}.
    This concerns the bound, and says nothing about $\sigma_L$.
    The present note derives $\sigma_L$ from the inherited lift and thereby converts the bound into a signal under a
    precisely named phase hypothesis.

    The discipline of the programme is respected throughout: no positive frontier run is read as a proof of $u\neq 0$.
    Every claim below carries an explicit status among \emph{proved}, \emph{structural}, \emph{conditional}, and
    \emph{open}.

  \section{Problem statement: $T(e)$ cannot define $\sigma_L$}

    The required parities are fixed by the oriented-frontier observable~\cite{Beau2026q11of}.
    Under the residual reflection $\varphi = R_b\circ\JPi$ one has $\dac\circ\varphi = -\dac$ and, under [H-WS]
    of~\cite{Beau2026q11of} (which gives $[\gamma_5,R_b]=0$), $\gamma_5\circ\varphi = -\gamma_5$, so the dressed signal $\gamma_5\dac$ is $\varphi$-even and survives the re-pairing that kills the raw
    cocycle.
    Under cascade reversal the cocycle is arrow-odd, $\dac(e^{-1}) = -\dac(e)$, so the reversal control
    $\Theta_\chi^{(\partial^+)^{-1}} = -\Theta_\chi^{\partial^+}$ holds only if the chiral weight is transported rather than
    recomputed.
    The admissible $\sigma_L$ must therefore, under [H-WS], satisfy
    \begin{equation}
      \sigma_L\circ\varphi = -\sigma_L \quad(\varphi\text{-odd}),
      \qquad
      \sigma_L(e^{-1}) = \sigma_L(e) \quad(\text{arrow-even}).
    \end{equation}
    By contrast $\sgn\dac$ is $\varphi$-odd and arrow-odd; the global identification $\sigma_L = \sgn\dac$ on all directed
    edges is therefore excluded, since it makes $\sigma_L$ arrow-odd and breaks the reversal control.

  \section{The $2\times 2$ obstruction}

    On the outer doublet $\{e_+,e_-\}$ of the $\JPi$-odd sector~\cite[Prop.~6.3]{Beau2026q14}, with the identifications
    \begin{equation}
      J_3 = \gamma_5 = \sigma_3,\qquad
      \JPi = \sigma_1 K \ (\text{antiunitary}),\qquad
      \Rmix = \begin{pmatrix}0 & 1\\ -1 & 0\end{pmatrix} = i\sigma_2,
    \end{equation}
    the $\JPi$-odd projection of the cascade generator is
    \begin{equation}
      T(e) = \alpha(e)\,\sigma_3 + \mu(e)\,i\sigma_2
      = \begin{pmatrix}\alpha & \mu\\ -\mu & -\alpha\end{pmatrix},
      \qquad \lambda_\pm = \pm\sqrt{\alpha^2-\mu^2},
    \end{equation}
    where $\alpha$ is the real-split coefficient produced by the order of the product through the $\mathfrak{sl}_2$
    commutator $[E,F]=H$ and $\mu$ is the mixing coefficient~\cite[Prop.~6.5]{Beau2026q14}.
    The conjugations $\JPi\sigma_3\JPi^{-1} = -\sigma_3$ and $\JPi(i\sigma_2)\JPi^{-1} = -i\sigma_2$ give $\JPi T
  \JPi^{-1} = -T$, so $T$ is $\JPi$-odd.

    \begin{lemma}[Parity collapse]
Let $U(e)$ be the lift of the oriented edge with $T(e) = \log U(e)$ restricted to the doublet.
Edge inversion $U(e^{-1}) = U(e)^{-1}$ gives $T(e^{-1}) = -T(e)$, and, under [H-WS] of~\cite{Beau2026q11of}, the
residual reflection gives $T(\varphi e) = \varphi T(e)\varphi^{-1} = -T(e)$.
Hence both operations act as $(\alpha,\mu)\mapsto(-\alpha,-\mu)$, and every function $f(T)$ has identical
arrow-parity and $\varphi$-parity.
    \end{lemma}

    \begin{proof}
      $U(e^{-1}) = U(e)^{-1} = \exp(-T(e))$ yields $T(e^{-1}) = -T(e)$.
      $T$ is $\JPi$-odd and, under [H-WS] of~\cite{Beau2026q11of}, $R_b$ acts on the chiral carrier as a scalar (part~(i); its value plays no role here), so $\varphi T\varphi^{-1} = -T$ and
      $T(\varphi e) = -T(e)$.
      A function $f(T)$ that is invariant (resp.\ anti-invariant) under $T\mapsto -T$ is therefore even (resp.\ odd) under
      both arrow inversion and $\varphi$ simultaneously.
    \end{proof}

    \begin{theorem}[$\sigma_L$ is not a function of $T(e)$]
There is no function $f$ such that $\sigma_L = f(T(e))$ realises the admissible parities, because the admissible
$\sigma_L$ is arrow-even and $\varphi$-odd, while every $f(T)$ has equal arrow- and $\varphi$-parities.
    \end{theorem}

    \begin{status}\textbf{Proved under [H-WS] of~\cite{Beau2026q11of}.}\end{status}

    The spectral branch confirms the obstruction explicitly.
    The projector onto the positive eigenline $v_+\propto(\mu,\ \lambda_+-\alpha)$ gives
    \begin{equation}
      \operatorname{Tr}(\gamma_5 P_{v_+})
      = \frac{\mu^2 - (\lambda_+-\alpha)^2}{\mu^2 + (\lambda_+-\alpha)^2}
      = \frac{\lambda_+}{\alpha}
      = \sgn(\alpha)\sqrt{1 - \mu^2/\alpha^2},
    \end{equation}
    whose sign is $\sgn(\alpha)$, hence arrow-odd for every value of $\mu$.
    Under the complexification $\mathrm{mp}(2,\mathbb{R})_{\mathbb{C}}\simeq\mathfrak{sl}_2(\mathbb{C})$ the real Cartan
    $J_3$ is the non-compact boost, so the spectral branch is exactly the boost sign, which is the excluded choice.

  \section{Representation-type resolution}

    The carrier $V$ is the Weyl spinor of $\mathfrak{sl}_2(\mathbb{C})$, with the fundamental $\two$ ($L$) and its complex
    conjugate $\twobar$ ($R$); the generation factor is $\mathrm{Sym}^2(V)$, so $e_+ = L\otimes L$ and $e_- = R\otimes
  R$~\cite[\S 6]{Beau2026q14}.
    The lift acts on $e_+$ by $U\otimes U$ and on $e_-$ by $\bar U\otimes\bar U$.
    The chiral weight is the representation type carried by the edge,
    \begin{equation}
      \sigma_L(e) = (\text{Weyl representation type of the lift of } e)\in\{+1\ (\two),\ -1\ (\twobar)\}.
    \end{equation}

    \begin{proposition}[Correct parities of the representation type]
Group inversion preserves the representation type, $\sigma_L(e^{-1}) = \sigma_L(e)$ (arrow-even), because $U^{-1}$ lies
in the same representation as $U$.
Under [H-WS] of~\cite{Beau2026q11of}, complex conjugation exchanges $\two\leftrightarrow\twobar$, $\sigma_L(\varphi e) =
-\sigma_L(e)$ ($\varphi$-odd), because $\varphi$ contains the antiunitary $\JPi$ and $R_b$ does not itself exchange the
two types.
    \end{proposition}

    \begin{status}\textbf{Structural, under [H-WS].}\end{status}

    This is consistent with the real symplectic model.
    For a real lift $U\in\mathrm{SL}(2,\mathbb{R})$ the restricted representation is self-conjugate: the distinction
    between $\two$ and $\twobar$ is not a canonical label on the real carrier.
    Thus the effective chiral contrast carried by this real model is zero, which we denote schematically by $\sigma_L = 0$
    before the complex Lorentzian lift is chosen.
    This matches the chiral symmetry of the real cascade, in which $\mu = 0$~\cite[Rem.~6.4]{Beau2026q14}.
    The representation type is the unique label distinguished by complex conjugation, which is precisely the operation
    separating $\varphi$ from arrow inversion; this is why it can carry opposite parities under the two.

  \section{The inherited central cyclotomic phase}

    In the Schr\"odinger--Heisenberg convention of the exact Weil-block construction~\cite{Beau2026a16}, with the group law
    $(a,b,\gamma)(a',b',\gamma')=(a+a',b+b',\gamma+\gamma'+ab')$, the lift acts as
    \begin{equation}
      \rho_c(a,b,\gamma)\,f(x) = \zq^{\,c(\gamma + bx)}\,f(x+a),
      \qquad \zq = e^{2\pi i/q},
    \end{equation}
    which satisfies $\rho_c(g)\rho_c(h)=\rho_c(gh)$ for this law,
    and has central character $\zq^{\,c\gamma}$.
    The central increment of an oriented edge is the Heisenberg cocycle, so the edge phase is the linear cyclotomic
    character
    \begin{equation}
      \omega_\chi(e) = \zq^{\,c\,a(g)\,s_b} = \zq^{\,\dac(e)}.
    \end{equation}
    The two chiral branches carry the conjugate central characters $c$ and $q-c$, which is exactly the parity
    $c\leftrightarrow q-c$ realised by $\JPi$~\cite[\S 6]{Beau2026q14}.

    \begin{proposition}[Arithmetic nature of the inherited phase]
In this convention the edge phase is purely linear in the cocycle: the exponent is the central increment $c\,a(g)\,s_b$ of the group law, with
no quadratic residue, and $|\omega_\chi(e)| = 1$ on every edge.
There is therefore no Gauss sum $\sqrt{q}$, no Maslov quadratic character, and no $\sqrt{5}$ amplitude in the
inherited phase.
    \end{proposition}

    \begin{status}\textbf{Proved within the inherited central cyclotomic phase convention.}
      The exponent $c\,a(g)\,s_b$ is the central increment of the group law and $|\zq^k|=1$ termwise; no separate
      numerical verification is claimed.\end{status}

    The absence of a Fourier or Weyl generator in the Schr\"odinger--Heisenberg pipeline is the structural reason no
    metaplectic Gauss phase appears: the lift is a product of Heisenberg translations and modulations, whose only phase is
    the inherited central character.

  \section{Frontier overlap and the conditional signal}

    On the nonzero-cocycle support of the outgoing frontier, the natural branch of the inherited central phase is selected
    by the sign of its imaginary part,
    \begin{equation}
      \sigma_L(e) = \sgn\,\Im\,\zq^{\,\dac(e)} = \sgn \dac(e),
      \qquad e\in\partial^+,
    \end{equation}
    since for a centred residue $r\in(-q/2,q/2)$ one has $\sgn\sin(2\pi r/q) = \sgn r$.
    This formula is an outgoing-frontier branch selection.
    For the paired inverted edge $e^{-1}$ the chiral label is transported with the same branch,
    \begin{equation}
      \sigma_L(e^{-1}) := \sigma_L(e),
    \end{equation}
    rather than recomputed from $\dac(e^{-1})$.
    This is the arrow-even transport required by the reversal control, realised as a label on the unoriented edge class
    with canonical radial representative, while complex conjugation still exchanges the two Weyl representation types and
    gives $\varphi$-oddness.
    The normalised overlap is then
    \begin{equation}
      \rhochi = \frac{\langle \sigma_L\,\dac\rangle_{\partial^+}}{\langle |\dac|\rangle_{\partial^+}} = +1,
    \end{equation}
    and, restricted to the cocycle support, the locking is maximal; this value is a consequence of the definition
    $\sigma_L:=\sgn\dac$ on the support, not a measured quantity.

    \begin{theorem}[Conditional orientation theorem]
In the inherited Heisenberg central-phase lift, the chiral weight $\sigma_L$ equals $\sgn\dac$ on the nonzero-cocycle
support of the outgoing frontier, transported arrow-even as a function of the unoriented edge.
Consequently, in the frontier normalisation used in this diagnostic,
\begin{equation}
  \omega_\chi(e) = \zq^{\,\dac(e)}
  \;\Longrightarrow\;
  \rhochi = 1
  \;\Longrightarrow\;
  \Horient
  \;\Longrightarrow\;
  \NA = \thmax \neq 0.
\end{equation}
The equality with $\thmax$ is a maximal-locking statement for the inherited central-phase frontier channel.
It is not an absolute calibration of the later spectral split.
    \end{theorem}

    \begin{status}\textbf{Proved within the inherited central cyclotomic phase convention and under [H-WS].}
      The overlap $\rhochi=1$ holds by the definition $\sigma_L=\sgn\dac$ on the support (it is not a measured
      quantity); the support fraction of the cocycle on $\partial^+S_m$ is
      $\tfrac13,\tfrac49,\tfrac{11}{27},\tfrac{24}{55},\tfrac{49}{109},\tfrac{84}{185}$ for $m=1,\dots,6$,
      independent of $q\ge 61$ (exact; \texttt{code/frontier\_exact.py}).\end{status}

    The reading is deliberately conservative: the exact frontier computation validates the backend and the bound; the analysis of the
    inherited lift, and only that, converts the bound into a signal.
    No step infers $\NA\neq 0$ from a positive run.

  \section{Residual caveat: the spin-Galois phase}

    The result above is the phase inherited from the Heisenberg central character, which is the phase carried by the
    radial capacity pipeline~\cite{Beau2026a16}.
    The full Lorentzian spin-lift, with a genuine $\mathrm{SL}(2,\mathbb{C})$ factor and the ADE-spin structure carrying
    $\sqrt{5}$, is not present in this pipeline.
    If the chiral sector carried an additional phase beyond the inherited central character, the dichotomy would reopen.

    \begin{remark}[The orthogonality alternative]
Should the chiral branch acquire a phase in a Galois factor orthogonal to $\zq$, in particular a $\sqrt{5}$ or
ADE-spin factor, then $\langle\sigma_L\,\sgn\dac\rangle_{\partial^+}$ could vanish by character
orthogonality~\cite{Beau2026aog}, yielding a clean no-go of the same type as the arithmetic obstruction governing the
ADE case selection of the generation gate, which is isolated in~\cite{Beau2026aog}.
In that case the signal would exist operationally while the frontier amplitude vanished.
    \end{remark}

    \begin{status}\textbf{Open.}
      Whether the chiral spin-lift carries only the inherited central character, or an additional spin-Galois phase, is the
      single remaining obstruction besides [H-WS]; it is analytical and lies outside the present frontier computation.\end{status}

  \section{Relation to the Schur transport}

    The conditional orientation theorem fixes the frontier amplitude datum $\NA$, the normalisation entering the angular
    generation split~\cite{Beau2026q11of,Beau2026bim}.
    It does not, by itself, establish the polarisation transport that opens the mass split.
    The split is read, through $\Phi$ of [H-Gen] in~\cite{Beau2026prs}, in the compression remainder of the projected
    Dirac residue, and its $\JPi$-odd component on the generation sector is
    \begin{equation}
      \PiS\, D\,(\partial_s P|_0)\, D\,\PiS^{*}\Big|_{\mathrm{gen},\,\JPi\text{-odd}} \neq 0,
    \end{equation}
    which is a distinct statement on the same chiral object~\cite{Beau2026prs} (its outer-weight part, read through
    $\Phi$ of [H-Gen] for the $\JPi$-odd part of $\partial_sP|_0$, is the Schur-transverse condition of that note).
    The two stages are sequential: $\Horient$ promotes the bound to a signal $\NA\neq 0$, and a transport of the
    oriented signal to the $\JPi$-odd generation sector would then promote it to the spectral split $u\neq 0$ within
    the stratigraphic mass hierarchy~\cite{Beau2026a4}.
    The only transport stated in~\cite{Beau2026prs} is the hypothesis [H-Tr] for the $\mathfrak{sl}_2$ opening
    $\alpha(x,y)$ along a family of locking operators, supplied by no source; the identification of the frontier
    signal with that opening is a further input.
    The Schur-null and Schur-transverse loci are defined there, and which one holds is open.

    \begin{status}\textbf{Not derived here; no transport from the frontier signal to the locking operator is supplied
      (cf.\ [H-Tr] of~\cite{Beau2026prs} for the $\mathfrak{sl}_2$ opening).}\end{status}

  \section{Status summary}

    \begin{itemize}
      \item $\sigma_L \neq f(T(e))$: \textbf{proved under [H-WS]}.
      \item $\sigma_L = $ Weyl representation type ($\two$ versus $\twobar$): \textbf{structural, under [H-WS]}.
      \item $\omega_\chi(e) = \zq^{\,\dac(e)}$ linear, no Gauss / Maslov / $\sqrt{5}$ component:
      \textbf{proved within the inherited central cyclotomic phase convention}.
      \item $\rhochi = 1$ on the cocycle support, hence $\NA = \thmax \neq 0$:
      \textbf{proved within the inherited central cyclotomic phase convention and under [H-WS]}.
      \item Absence of an additional spin-Galois phase: \textbf{open}.
      \item $u \neq 0$: \textbf{not derived here; no transport from the frontier signal to the locking operator is supplied}.
    \end{itemize}

    The qualitative segment of $\Horient$ is thus closed in the inherited lift under [H-WS], with the residual obstruction
    reduced, besides [H-WS], to a single, sharply named analytical question: the spin-Galois phase of the full Lorentzian
    chiral lift.

    \newpage
    \bibliographystyle{plain}
    \bibliography{cosmochrony-bibliography}

\end{document}
